%%% FIELD proof-long-history-frontier
Fix integers \(n,T\geq 1\) and \(c\in[0,1]\) such that \(T\geq 2k_n\).  First consider the upper bounds.  For observations drawn under any \(P\in\mathcal P_T\), define
\[
\overline Y_n=\frac1n\sum_{i=1}^nY_i,
\qquad
\widehat\psi_n=\operatorname{proj}_{[0,1]}(\overline Y_n).
\tag{1}
\]
By \(\mathrm{ass:bounded\text{-}outcome}\), \(\overline Y_n\in[0,1]\), so the projection in (1) is in fact the identity; it makes the estimator a total element of \(\mathcal E_{n,T}\).  Put \(m=\mu(P)=E_PY\).  By \(\mathrm{ass:iid\text{-}sampling}\), \(E_{P^n}\overline Y_n=m\) and \(\operatorname{Var}_{P^n}(\overline Y_n)=\operatorname{Var}_P(Y)/n\).  Since \(0\leq Y\leq1\), one has \(Y^2\leq Y\), and hence
\[
\operatorname{Var}_P(Y)=E_P[Y^2]-m^2
\leq m-m^2\leq\frac14.
\tag{2}
\]
The cross term in the following bias--variance decomposition vanishes because \(E_{P^n}(\overline Y_n-m)=0\).  Thus (2) and \(\mathrm{lem:stable\text{-}weights}\) give
\[
\begin{aligned}
E_{P^n}\!\left[(\widehat\psi_n-\psi_T(P,c))^2\right]
&=\frac{\operatorname{Var}_P(Y)}n+\{m-\psi_T(P,c)\}^2\\
&\leq\frac1{4n}+B(c)^2
\leq\frac1{4n}+\frac{c^2}{9}
\leq n^{-1}+c^2.
\end{aligned}
\tag{3}
\]
Both \(\widehat\psi_n\) and \(\psi_T(P,c)\) belong to \([0,1]\), so the left side of (3) is also at most one.  Taking the supremum over \(P\in\mathcal P_T\) and then using \(\mathrm{def:minimax\text{-}risk}\) yields
\[
R(n,T,c)\leq\min\{1,n^{-1}+c^2\}.
\tag{4}
\]

We next construct the advertised interval without invoking an asymptotic approximation.  Let
\[
h_n=\sqrt{\frac{\log(40)}{2n}},
\qquad
I_n=\bigl[\overline Y_n-h_n-B(c),\,\overline Y_n+h_n+B(c)\bigr]\cap[0,1].
\tag{5}
\]
This is a measurable connected interval and is defined for every possible sample.  To verify its coverage from first principles, convexity of \(y\mapsto e^{\lambda y}\) on \([0,1]\) gives
\[
e^{\lambda y}\leq1-y+ye^\lambda
\qquad(0\leq y\leq1,\ \lambda\in\mathbb R).
\tag{6}
\]
Taking expectations in (6), the centered log-moment-generating function is bounded by
\[
\log E_Pe^{\lambda(Y-m)}
\leq g_m(\lambda)
:=\log(1-m+me^\lambda)-m\lambda.
\tag{7}
\]
Writing \(p_\lambda=me^\lambda/(1-m+me^\lambda)\), direct differentiation gives
\[
g_m(0)=g_m'(0)=0,
\qquad
g_m''(\lambda)=p_\lambda(1-p_\lambda)\leq\frac14.
\tag{8}
\]
Twice integrating (8) from zero to \(\lambda\) shows that \(g_m(\lambda)\leq\lambda^2/8\) for every real \(\lambda\).  Independence in \(\mathrm{ass:iid\text{-}sampling}\), exponential Markov inequality, and optimization at \(\lambda=4h_n\) therefore give
\[
P^n\{\overline Y_n-m>h_n\}
\leq\inf_{\lambda>0}\exp\{-n\lambda h_n+n\lambda^2/8\}
=e^{-2nh_n^2}.
\tag{9}
\]
Applying the same argument to \(m-\overline Y_n\), and then the union bound, gives
\[
P^n\{|\overline Y_n-m|>h_n\}
\leq2e^{-2nh_n^2}=0.05.
\tag{10}
\]
On the complement of the event in (10), \(\mathrm{lem:stable\text{-}weights}\) implies
\[
|\overline Y_n-\psi_T(P,c)|
\leq|\overline Y_n-m|+|m-\psi_T(P,c)|
\leq h_n+B(c).
\tag{11}
\]
Because \(\psi_T(P,c)\in[0,1]\), (11) implies \(\psi_T(P,c)\in I_n\).  Thus (10) proves that \(I_n\in\mathcal C_{n,T,c}\).  Moreover,
\[
\begin{aligned}
\operatorname{length}(I_n)
&\leq\min\{1,2h_n+2B(c)\}\\
&\leq\min\!\left\{1,C_L(n^{-1/2}+c)\right\}
\leq C_L\min\{1,n^{-1/2}+c\},
\end{aligned}
\tag{12}
\]
where \(C_L=2\sqrt{\log(40)/2}\geq2/3\).  Taking the supremum over \(P\) and using \(\mathrm{def:minimax\text{-}length}\) proves the desired length upper bound.  Computing (1) and (5) requires \(O(n)\) arithmetic operations and \(O(1)\) additional memory.

It remains to prove the two converses.  Since \(k_n\geq1\), the assumed inequality \(T\geq2k_n\) implies \(T>k_n\) and
\[
1-\frac{k_n}{T}\geq\frac12.
\tag{13}
\]
The final two lower bounds of \(\mathrm{prop:finite\text{-}history\text{-}envelopes}\), which already take the infimum over all measurable estimators and over all possibly randomized honest connected intervals in the original iid experiment, therefore imply for its universal constant \(a_1>0\) that
\[
\begin{aligned}
R(n,T,c)
&\geq a_1\left\{n^{-1}+c^2\left(1-\frac{k_n}{T}\right)\right\}
\geq\frac{a_1}{2}(n^{-1}+c^2),\\
L(n,T,c)
&\geq a_1\left\{n^{-1/2}+c\sqrt{1-\frac{k_n}{T}}\right\}
\geq\frac{a_1}{\sqrt2}(n^{-1/2}+c).
\end{aligned}
\tag{14}
\]
Each last expression in (14) is at least the same constant times the corresponding quantity truncated at one.  Consequently, with
\[
a_0=\min\{a_1/2,1\},
\qquad
A_0=\max\{1,C_L\},
\tag{15}
\]
equations (4), (12), and (14) prove all four asserted inequalities, with \(0<a_0\leq A_0<\infty\).  None of these constants depends on \(n,T,c\).  The argument is uniform over the whole class \(\mathcal P_T\), so it includes shrinking \(c\), deterministic treatment kernels, and null histories; no division by a propensity was introduced.
